% Revisión: permite cortar las URL largas de la bibliografía también en los guiones
\PassOptionsToPackage{hyphens}{url}
\usepackage[utf8]{inputenc}
\usepackage[T1]{fontenc}
\usepackage[spanish,es-tabla]{babel}

\decimalpoint

\usepackage[a4paper]{geometry} % Tamaño del papel
\usepackage{graphicx}               % Para incluir imágenes
\usepackage{subcaption}
\usepackage{amsmath}                % Para el manejo de matemáticas
\usepackage{amsthm}
\usepackage{multirow}					% Combinar columnas
\usepackage{url}					% Dar formato a urls				
\usepackage{array}					% Para ajustar el texto en la celda
\usepackage[table]{xcolor}
\usepackage{hhline}					% Líneas horizontales que no tapan las celdas de color
\usepackage{longtable}				% Tablas que se dividen entre páginas (tablas de riesgos)
% Colores pastel de los niveles de riesgo (matriz y tablas de riesgos); el medio es el mismo amarillo de los diagramas
\definecolor{riskLow}{RGB}{212,237,218}
\definecolor{riskMedium}{RGB}{249,232,160}
\definecolor{riskHigh}{RGB}{248,203,173}
\definecolor{riskVeryHigh}{RGB}{240,160,160}
% Celda que ocupa las 5 columnas de la derecha en las tablas de riesgos (análisis de riesgos)
\newcommand{\riskSpan}[1]{\multicolumn{5}{p{\dimexpr 0.79\textwidth+8\tabcolsep+4\arrayrulewidth\relax}|}{#1}}
\usepackage{float}
\usepackage{booktabs}
\usepackage{comment} 
\usepackage{hyperref}
\usepackage{enumitem}
\usepackage{amssymb}		% Simbolos matematicos
\usepackage{fancyhdr}  % Para encabezados personalizados
\usepackage{algorithm}
\usepackage{algpseudocode}
\usepackage{verbatim}
\usepackage[edges]{forest}
\usepackage{tabularx}
\usepackage{tikz}
\usetikzlibrary{positioning, arrows.meta, shadows}

% Estilo de los diagramas de flujo (colores del diseño original en Mermaid)
\definecolor{boxFill}{RGB}{249,232,160}
\definecolor{boxBorder}{RGB}{225,202,114}
\definecolor{boxText}{RGB}{40,37,61}
\definecolor{labelFill}{RGB}{214,214,214}
\tikzset{
	flowchart/.style={node distance=2.2cm, font=\sffamily},
	flowbox/.style={draw=boxBorder, fill=boxFill, text=boxText, line width=0.8pt,
		rounded corners=3pt, drop shadow={opacity=0.18, shadow xshift=0.6pt, shadow yshift=-0.8pt},
		minimum width=2.6cm, minimum height=1cm, inner xsep=8pt},
	flowarrow/.style={-{Stealth[length=7pt, width=6pt]}, line width=1pt, rounded corners=12pt},
	flowlabel/.style={font=\sffamily\small, fill=labelFill, text=boxText,
		inner sep=2.5pt, rounded corners=0pt},
}

% Revisión: pies de figura y tabla más angostos que el texto (no aplica a las subfiguras)
\usepackage{caption}
\captionsetup{margin=0.075\textwidth}
\captionsetup[sub]{margin=0pt}
% Revisión: ninguna figura o tabla aparece antes de su primera referencia (flafter) ni fuera de la
% sección, subsección o subsubsección en la que se cita (placeins)
\usepackage{flafter}
\usepackage[section]{placeins}
\usepackage{etoolbox}
\pretocmd{\subsection}{\FloatBarrier}{}{}
\pretocmd{\subsubsection}{\FloatBarrier}{}{}

% Listas compactas dentro de tablas: sin sangría, sin renglón vacío arriba y sin ancla de hyperref
% (el ancla empuja la lista un renglón hacia abajo). Solo se activa en las especificaciones de casos de uso
\newif\ifcompactTableLists
\makeatletter
\AtBeginEnvironment{tabular}{\ifcompactTableLists\setlist[itemize,enumerate]{nosep, wide, labelindent=0pt,
	before={\@minipagetrue\let\refstepcounter\stepcounter}}\fi}
\makeatother

% Algoritmo que no flota y puede partirse entre páginas (mismo estilo "ruled" del paquete algorithm).
% Se arma en una caja; si no cabe en lo que queda de la página se corta con \vsplit, la parte de abajo
% se cierra con una línea y la página siguiente empieza con "Algorithm N (continuación)".
\makeatletter
\newsavebox\algorithmBox
\newdimen\algorithmRoom
\newif\ifalgorithmStarted
\newenvironment{breakablealgorithm}
	{\par\addvspace{\medskipamount}%
	\refstepcounter{algorithm}%
	\setbox\algorithmBox\vbox\bgroup
		\hsize\linewidth\@parboxrestore
		\hrule height.8pt depth0pt\kern2pt%
		\renewcommand{\caption}[2][\relax]{%
			{\raggedright\textbf{\fname@algorithm~\thealgorithm} ##2\par}%
			\ifx\relax##1\relax
				\addcontentsline{loa}{algorithm}{\protect\numberline{\thealgorithm}##2}%
			\else
				\addcontentsline{loa}{algorithm}{\protect\numberline{\thealgorithm}##1}%
			\fi
			\kern2pt\hrule\kern2pt}%
	}
	{\par\egroup
	\algorithmStartedfalse
	\begingroup
		\splittopskip=\ht\strutbox \splitmaxdepth=\maxdepth \vbadness=10000
		\breakablealgorithm@loop
	\endgroup
	\par\addvspace{\medskipamount}}
% Espacio que queda en la página, reservando la línea de cierre
\def\breakablealgorithm@room{%
	\ifdim\pagegoal=\maxdimen \algorithmRoom=\textheight \else
		\algorithmRoom=\pagegoal \advance\algorithmRoom by -\pagetotal \fi
	\advance\algorithmRoom by -\baselineskip}
\def\breakablealgorithm@loop{%
	\breakablealgorithm@room
	\dimen@=\ht\algorithmBox \advance\dimen@ by \dp\algorithmBox
	\ifdim\dimen@>\algorithmRoom
		\ifdim\algorithmRoom<4\baselineskip
			% No cabe ni el inicio: se pasa todo a la página siguiente
			\newpage
			\ifalgorithmStarted\breakablealgorithm@header\fi
		\else
			\setbox\z@=\vsplit\algorithmBox to \algorithmRoom
			\setbox\z@=\vbox{\unvbox\z@}%
			\box\z@ \nointerlineskip\kern2pt\hrule
			\algorithmStartedtrue
			\newpage
			\breakablealgorithm@header
		\fi
		\expandafter\breakablealgorithm@loop
	\else
		\box\algorithmBox \nointerlineskip\kern2pt\hrule
	\fi}
\def\breakablealgorithm@header{%
	\hrule height.8pt depth0pt\kern2pt
	{\raggedright\textbf{\fname@algorithm~\thealgorithm} (continuación)\par}%
	\kern2pt\hrule\kern2pt}
\makeatother

% \bstctlcite de IEEEtran (la clase report no lo define): pasa opciones de control a IEEEtran.bst
\makeatletter
\def\bstctlcite{\@ifnextchar[{\@bstctlcite}{\@bstctlcite[@auxout]}}
\def\@bstctlcite[#1]#2{\@bsphack
	\@for\@citeb:=#2\do{%
		\edef\@citeb{\expandafter\@firstofone\@citeb}%
		\if@filesw\immediate\write\csname #1\endcsname{\string\citation{\@citeb}}\fi}%
	\@esphack}
\makeatother

\newtheorem{definition}{Definición}
\setlength{\parskip}{10pt}

% Pendientes visibles en el PDF: \TODO{texto}
\newcommand{\TODO}[1]{\par\noindent\fcolorbox{orange!70!black}{yellow!25}{\parbox{\dimexpr\linewidth-2\fboxsep-2\fboxrule\relax}{\textbf{TODO:} #1}}\par}
% Revisión: tercera pasada de TeX solo en los párrafos que de otro modo se salen del margen
\setlength{\emergencystretch}{3em}

% Configurar encabezados
\pagestyle{fancy}
\fancyhf{}
\fancyhead[LE,RO]{\small\leftmark}    % Muestra el capítulo en lugar de la sección
\fancyhead[RE,LO]{\small\thepage}      % Número de página
\setlength{\headheight}{13.6pt}

\setcounter{secnumdepth}{4}
\setcounter{tocdepth}{4}

%%% Botones - BTN - Buttons
\newcounter{BTNcounter}
\newcommand{\NewBTN}[2]{%
	\refstepcounter{BTNcounter}%
	\expandafter\gdef\csname BTNtext@#1\endcsname{#2}%
	\hypertarget{#1}{}%  Define el destino del enlace
\label{#1}%
	\textcolor{black}{BTN\theBTNcounter~\csname BTNtext@#1\endcsname}%
}
\newcommand{\BTN}[1]{\hyperlink{#1}{\textcolor{black}{BTN\ref{#1}~\csname BTNtext@#1\endcsname}}}

\newcounter{SLDcounter}
\newcommand{\NewSLD}[2]{%
	\refstepcounter{SLDcounter}%
	\expandafter\gdef\csname SLDtext@#1\endcsname{#2}%
	\hypertarget{#1}{}%  Define el destino del enlace
	\label{#1}%
	\textcolor{black}{SLD\theSLDcounter~\csname SLDtext@#1\endcsname}%
}
\newcommand{\SLD}[1]{\hyperlink{#1}{\textcolor{black}{SLD\ref{#1}~\csname SLDtext@#1\endcsname}}}

%%% Inputs - IN 
\newcounter{INcounter}
\newcommand{\NewIN}[2]{%
	\refstepcounter{INcounter}%
	\expandafter\gdef\csname INtext@#1\endcsname{#2}%
	\hypertarget{#1}{}%  Define el destino del enlace
\label{#1}%
	\textcolor{black}{IN\theINcounter~\csname INtext@#1\endcsname}%
}
\newcommand{\IN}[1]{\hyperlink{#1}{\textcolor{black}{IN\ref{#1}~\csname INtext@#1\endcsname}}}

%%% Outputs - OUT 
\newcounter{OUTcounter}
\newcommand{\NewOUT}[2]{%
	\refstepcounter{OUTcounter}%
	\expandafter\gdef\csname OUTtext@#1\endcsname{#2}%
	\hypertarget{#1}{}%  Define el destino del enlace
\label{#1}%
	\textcolor{black}{OUT\theOUTcounter~\csname OUTtext@#1\endcsname}%
}
\newcommand{\OUT}[1]{\hyperlink{#1}{\textcolor{black}{OUT\ref{#1}~\csname OUTtext@#1\endcsname}}}

%%% Messages - MSG
\newcounter{MSGcounter}
\newcommand{\NewMSG}[2]{%
	\refstepcounter{MSGcounter}%
	\expandafter\gdef\csname MSGtext@#1\endcsname{#2}%
	\hypertarget{#1}{}%  Define el destino del enlace
	\label{#1}%
	\textcolor{black}{MSG\theMSGcounter~\csname MSGtext@#1\endcsname}%
}
\newcommand{\MSG}[1]{\hyperlink{#1}{\textcolor{black}{MSG\ref{#1}~\csname MSGtext@#1\endcsname}}}

%%% User Interface - UI
\newcounter{UIcounter}
\newcommand{\NewUI}[2]{%
	\refstepcounter{UIcounter}%
	\expandafter\gdef\csname UItext@#1\endcsname{#2}%
	\hypertarget{#1}{}%  Define el destino del enlace
\label{#1}%
	\textcolor{black}{UI\theUIcounter~\csname UItext@#1\endcsname}%
}
\newcommand{\UI}[1]{\hyperlink{#1}{\textcolor{black}{UI\ref{#1}~\csname UItext@#1\endcsname}}}

\setlength{\headheight}{27.1pt} 
\useforestlibrary{edges}
\newcommand{\fitem}[2]{\makebox[3cm][l]{#1}{\textcolor{gray}{#2}}}
